% !PRINTED none
\begin{center}
\gesperrt{Sixteenth Chapter.}

\vspace{1em}

\textbf{Class Number in Quadratic Fields.}

\vspace{0.5em}
\rule{4cm}{0.4pt}
\end{center}

\begin{center}
§ 110. \gesperrt{Fundamental Units in the Orders.}
\end{center}

The theory of units, which we have set forth in general in Bd. II,
§ 191, takes on a simpler shape for the quadratic fields, but must on
the other hand be extended and elaborated for the \gesperrt{orders}.

The integers of the order $O$ with the discriminant $D$ are, by
§ 96, (6), contained in the form $x_1 + x_2\theta$, where
\[
\begin{array}{ll}
\theta = \tfrac{1}{2}\sqrt{D}, & \text{if } D \equiv 0 \ (\mathrm{mod}\ 4),\\[4pt]
\theta = \dfrac{1+\sqrt{D}}{2}, & \text{„}\quad D \equiv 1 \ (\mathrm{mod}\ 4),
\end{array}
\]
and $x_1, x_2$ are integral rational numbers.

Such a number $\varepsilon$ is a unit if

\weq{1}{N(\varepsilon) = (x_1+x_2\theta)(x_1+x_2\theta') = \pm 1}

holds. From this there results the condition

\weq{2}{x_1^2 - x_2^2 D = \pm 4.}

If $D$ is negative, then only the positive sign can stand here, and in
general there are only the two solutions $x = \pm 2$,
$y = 0$ % [sic] German source prints x and y instead of x_1 and x_2., and in the two exceptional cases the four or six
solutions:
\[
\begin{array}{lll}
D = -4, & x_1 = \pm 2, & x_2 = 0,\\
 & x_1 = 0, & x_2 = \pm 1.\\
D = -3, & x_1 = \pm 2, & x_2 = 0,\\
 & x_1 = \pm 1, & x_2 = \pm 1.
\end{array}
\]
